\exercisesection{Normed Algebras}

\begin{enumerate}
\def\labelenumi{\arabic{enumi})}
\tightlist
\item
  Let (A, e) be a unital complex algebra.
\end{enumerate}

\begin{enumerate}
\def\labelenumi{\alph{enumi})}
\tightlist
\item
  For a linear form \(\lambda \in \mathrm{A}^*\) such that \(\lambda(e) = 1\) , the following conditions are equivalent:
\end{enumerate}

\begin{enumerate}
\def\labelenumi{(\roman{enumi})}
\item
  For every \(x \in \operatorname{Ker}(\lambda)\) , one has \(x^2 \in \operatorname{Ker}(\lambda)\) ;
\item
  For every \(x \in \mathbb{A}\) , one has \(\lambda(x^{2}) = \lambda(x)^{2}\) ;
\item
  The kernel of \(\lambda\) is a right ideal of A;
\item
  The linear form \(\lambda\) is a unital morphism from A to C.
\end{enumerate}

(Observe that \(\lambda(x - \lambda(x)e) = 0\) for all \(x\in A\) .)

\begin{enumerate}
\def\labelenumi{\alph{enumi})}
\setcounter{enumi}{1}
\item
  Let \(\lambda \in \mathrm{A}^*\) a linear form such that \(\lambda(1) = 1\) and such that for every \(x \in \mathrm{A}\) , one has \(\lambda(x) \in \mathrm{Sp}_{\mathrm{A}}(x)\) . Let \(x \in \mathrm{Ker}(\lambda)\) such that \(\lambda(x^{2}) \neq 0\) The spectrum of \(x\) is not bounded in \(\mathbf{C}\) . (For \(n \geqslant 2\) integer, consider the polynomial \(p \in \mathbf{C}[\mathrm{X}]\) such that \(p(z) = \lambda((ze - x)^{n})\) for all \(z \in \mathbf{C}\) ; observe that its roots belong to the spectrum of \(x\) and bound their Euclidean norm below in terms of \(|\lambda(x^{2})|\) .)
\item
  Suppose that A is a Banach algebra. For a linear form \(\lambda\) on A to be a unital morphism from A into C, it is necessary and sufficient that \(\lambda(e)=1\) and that \(\lambda(x)\in\mathrm{Sp}_{\mathrm{A}}(x)\) for all \(x\in A\) .
\end{enumerate}

\begin{enumerate}
\def\labelenumi{\arabic{enumi})}
\setcounter{enumi}{1}
\tightlist
\item
  Let \(n \geqslant 0\) be an integer. Let \(A_n\) the commutative complex Banach algebra of functions \(f: [0,1] \to \mathbf{C}\) admitting continuous derivatives in \([0,1]\) up to order \(n\) , equipped with the norm
\end{enumerate}

\[
\| f \| = \sum_ {k = 0} ^ {n} \frac {1}{k !} \sup _ {0 \leqslant t \leqslant 1} | f ^ {(k)} (t) |
\]

(example 4 of I, p.~18). Show that \(A_{n}\) is isomorphic to \(A_{m}\) if, and only if, n = m. (However, according to Example 1 of I, p.~144, these Banach algebras all have the same character space.)

\begin{enumerate}
\def\labelenumi{\arabic{enumi})}
\setcounter{enumi}{2}
\tightlist
\item
  Let G be a group equal to its derived subgroup (A, I, p.~67). For every \(g \in G\) , we denote by \(\ell(g)\) the smallest integer \(k \geqslant 0\) such that there exist pairs \((a_{1}, b_{1}), \ldots, (a_{k}, b_{k})\) into \(G \times G\) satisfying \(g = [a_{1}, b_{1}] \cdots [a_{k}, b_{k}]\) .
\end{enumerate}

\begin{enumerate}
\def\labelenumi{\alph{enumi})}
\item
  The map \(\ell\) is a metric on G.
\item
  For every \(g \in \mathbf{G}\) , the limit
\end{enumerate}

\[
\operatorname{lsc} (g) = \lim _ {n \to + \infty} \frac {\ell (g ^ {n})}{n}
\]

exists. It is called the stable commutator length of g.

\begin{enumerate}
\def\labelenumi{\alph{enumi})}
\setcounter{enumi}{2}
\tightlist
\item
  Let S be a set and \(\mathrm{G} = \mathrm{D}(\mathrm{F}(\mathrm{S}))\) the derived subgroup of the free group generated by S. The group G is equal to its derived subgroup, and for every \(g \in G\) , \(g \neq e\) , one has \(\operatorname{lsc}(g) \geqslant 1/2\) . There exists \(g \in G\) such that \(\operatorname{lsc}(g) = 1/2\) .
\end{enumerate}

\begin{enumerate}
\def\labelenumi{\arabic{enumi})}
\setcounter{enumi}{3}
\tightlist
\item
  Let \(n \geqslant 1\) an integer. We denote \(S(n)\) the largest integer \(k \geqslant 0\) such that, for every set \(X\) of \(n\) real numbers, there exists a subset \(Y \subset X\) of cardinality \(k\) such that the equation \(x + y = z\) has no solution in \(Y\) . We therefore have \(S(n) \leqslant n\) .
\end{enumerate}

\begin{enumerate}
\def\labelenumi{\alph{enumi})}
\tightlist
\item
  Show that the limit
\end{enumerate}

\[
\sigma = \lim _ {n \to + \infty} \frac {\mathrm{S} (n)}{n}
\]

exists.

\begin{enumerate}
\def\labelenumi{\alph{enumi})}
\setcounter{enumi}{1}
\item
  Show that \(\sigma \leqslant 2 / 5\) .
\item
  Show that \(\sigma \geqslant 1 / 3\) .
\end{enumerate}

(One can show that \(\sigma = 1/3\) , cf.~S. EBERHARD, B. GREEN, F. MANNERS, Sets of integers with no large sum-free subset, Annals of Mathematics 180 (2014), 621--652.)

\begin{enumerate}
\def\labelenumi{\arabic{enumi})}
\setcounter{enumi}{4}
\tightlist
\item
  Let A be a commutative ring and \(n \geqslant 1\) an integer. For a matrix \(M = (m_{i,j})\) of type \((n, n)\) with coefficients in A, we call the permanent of M the element
\end{enumerate}

\[
\mathrm{per(M)} = \sum_ {\sigma \in \mathrm{S} _ {n}} m _ {1, \sigma (1)} \dots m _ {n, \sigma (n)}
\]

of A.

Let \(k \geqslant 1\) an integer. We denote \(X_{n,k}\) the set of matrices M of type \((n, n)\) with real coefficients such that each column of M contains k coefficients equal to 1, all the others being zero.

Show that the limit

\[
\lim _ {n \to + \infty} \left(\inf _ {\mathrm{M} \in \mathrm{X} _ {n, k}} \operatorname{per} (\mathrm{M})\right) ^ {1 / n}
\]

exists.

\begin{enumerate}
\def\labelenumi{\arabic{enumi})}
\setcounter{enumi}{5}
\tightlist
\item
  Let \(k \geqslant 1\) an integer. For every integer \(N \geqslant 1\) , we denote by \(S_k(N)\) the largest cardinality of a set \(A \subset \{1, \ldots, N\}\) that does not contain \(k\) elements in arithmetic progression, that is, for any \(n_0 \geqslant 0\) , \(h \geqslant 1\) , \(\ell \geqslant 1\) , if \(\{n_0 + h, n_0 + 2h, \ldots, n_0 + \ell h\} \subset A\) , then \(\ell < k\) .
\end{enumerate}

\begin{enumerate}
\def\labelenumi{\alph{enumi})}
\tightlist
\item
  Show that the limit
\end{enumerate}

\[
s _ {k} = \lim _ {\mathrm{N} \to + \infty} \frac {\mathrm{S} _ {k} (\mathrm{N})}{\mathrm{N}}
\]

exists.

\begin{enumerate}
\def\labelenumi{\alph{enumi})}
\setcounter{enumi}{1}
\tightlist
\item
  Show that \(s_{1}=s_{2}=0\) One can show (“Szemerédi's theorem”, cf.~E. SZEMERÉDI, On sets of integers containing no k elements in arithmetic progression, Acta Arithmetica XXVII, 1975, 199--245) that \(s_{k}=0\) for all \(k\geqslant1\) ; for the case k=3, cf.~exercise 23 of II, p.~270.
\end{enumerate}

\begin{enumerate}
\def\labelenumi{\arabic{enumi})}
\setcounter{enumi}{6}
\tightlist
\item
  Let A be a normed algebra.
\end{enumerate}

\begin{enumerate}
\def\labelenumi{\alph{enumi})}
\tightlist
\item
  Let S be a nonempty bounded subset of A. Show that the limit
\end{enumerate}

\[
\varrho (\mathrm{S}) = \lim _ {n \to + \infty} \sup _ {x \in \mathrm{S} ^ {n}} \| x \| ^ {1 / n}
\]

exists and that \(\varrho(\mathrm{S}) = \varrho(x)\) if S is reduced to a single element x. We say that \(\varrho(\mathrm{S})\) is the spectral radius of the part S.

\begin{enumerate}
\def\labelenumi{\alph{enumi})}
\setcounter{enumi}{1}
\tightlist
\item
  Let N be the set of norms p on A, equivalent to the norm of A, and such that A, equipped with p, is a normed algebra. Let X be a subset of A. For there to exist \(p \in N\) such that \(p(x) \leqslant 1\) for all \(x \in X\) , it is necessary and sufficient that there exist a real number \(c \geqslant 0\) such that
\end{enumerate}

\[
\sup _ {n \geqslant 1} \sup _ {x \in X ^ {n}} \| x \| \leqslant c.
\]

\begin{enumerate}
\def\labelenumi{\alph{enumi})}
\setcounter{enumi}{2}
\tightlist
\item
  Show that
\end{enumerate}

\[
\varrho (\mathrm{S}) = \inf _ {p \in \mathrm{N}} \sup _ {x \in \mathrm{S}} p (x).
\]

\begin{enumerate}
\def\labelenumi{\arabic{enumi})}
\setcounter{enumi}{7}
\item
  Let A be the algebra of continuous complex-valued functions on \(\mathbf{R}\) tending to 0 at infinity, equipped with the norm \(\|f\| = \sup_{t\in\mathbf{R}} |f(t)|\) . Then \(\widetilde{\mathrm{A}}\) is identified with the algebra of continuous complex-valued functions on \(\mathbf{R}\) tending toward a finite limit at infinity. On \(\widetilde{\mathrm{A}}\) , the norm \(\|g\| = \sup_{t\in\mathbf{R}} |g(t)|\) differs from the norm defined in no. 1 of I, p.~15.
\item
  Let A be a normed algebra and \(b = \inf_{x \neq 0} \|x^{2}\|/\|x\|^{2}\) . Then
\end{enumerate}

\[
b \leqslant \inf _ {x \neq 0} \frac {\varrho (x)}{\| x \|} \leqslant \sqrt {b}.
\]

\begin{enumerate}
\def\labelenumi{\arabic{enumi})}
\setcounter{enumi}{9}
\item
  Let A be a normed algebra and \(x, y \in A\) . We have \(\varrho(xy) = \varrho(yx)\) .
\item
  Let H be a Hilbert space of countable type and \((f_{n})_{n\geqslant1}\) an orthonormal basis of H. Let A be the normed algebra \(\mathcal{L}(\mathrm{H})\) . Let \(x\in A\) the element defined by \(x(f_{i})=2^{-i}f_{i+1}\) . Show that x is quasi-nilpotent, but not nilpotent.
\end{enumerate}

\begin{enumerate}
\def\labelenumi{¶ \arabic{enumi})}
\setcounter{enumi}{11}
\tightlist
\item
  Let H be a separable Hilbert space and let \((f_n)_{n \geqslant 1}\) be an orthonormal basis of H. One defines the numbers \((\alpha_m)_{m \geqslant 1}\) by \(\alpha_m = e^{-k}\) if \(m\) is the product of \(2^k\) by an odd number.
\end{enumerate}

\begin{enumerate}
\def\labelenumi{\alph{enumi})}
\tightlist
\item
  Let \(x \in \mathcal{L}(\mathrm{H})\) defined by \(x(f_m) = \alpha_m f_{m+1}\) for \(m \geqslant 1\) Prove that \(x\) is not quasi-nilpotent. (Observe that
\end{enumerate}

\[
\left\| x ^ {n} \right\| = \sup _ {m} \bigl (\alpha_ {m} \alpha_ {m + 1} \dots \alpha_ {m + n - 1} \bigr),
\]

and evaluate \(\alpha_{1}\alpha_{2}\ldots\alpha_{2^{t}-1}\) .)

\begin{enumerate}
\def\labelenumi{\alph{enumi})}
\setcounter{enumi}{1}
\tightlist
\item
  Let \(x_{k} \in \mathcal{L}(\mathrm{H})\) defined by \(x_{k}(f_{m}) = 0\) if m is the product of \(2^{k}\) by an odd number, and \(x_{k}(f_{m}) = \alpha_{m} f_{m+1}\) otherwise. Show that \(x_{k}\) is nilpotent and that \((x_{k})\) converges to x as k tends to \(+\infty\) .
\end{enumerate}

\begin{enumerate}
\def\labelenumi{\arabic{enumi})}
\setcounter{enumi}{12}
\item
  Let A be a unital Banach algebra. Let \(x \in A\) such that \(\|x - 1\| < 1\) . There exists \(y \in A\) such that \(y^{2} = x\) .
\item
  Let A be a unital complex normed algebra. If \(\|x^{-1}\| = \|x\|^{-1}\) for every invertible element x of A, one has A = C · 1. (Reduce to the case where A is complete. Let \(x \in A\) , at a distance \(\alpha > 0\) from C · 1. If \(\lambda_{0} \in C\) is such that \(x - \lambda_{0}\) were invertible, then \(x - \lambda\) would be invertible for \(|\lambda - \lambda_{0}| < \alpha\) . Deduce that the spectrum of x would be empty.)
\item
  Let A be a unital complex normed algebra in which the only left topological divisor of zero is 0 and the only right topological divisor of zero is 0. Then \(A = C \cdot 1\) . (If \(x \in A\) , and if \(\lambda\) belongs to the boundary of \(\mathrm{Sp}_{\widehat{\mathrm{A}}}(\boldsymbol{x})\) , then \(x - \lambda\) is a topological divisor of zero in \(\widehat{A}\) , hence in A, therefore \(x = \lambda\) .)
\item
  Let A be a unital complex Banach algebra. Let \(a\) an element of A. The singular spectrum of A is called the set of \(\lambda \in \mathbf{C}\) such that \(a - \lambda e\) is a topological zero divisor on the left and on the right in A.
\end{enumerate}

\begin{enumerate}
\def\labelenumi{\alph{enumi})}
\item
  Show that the singular spectrum of A is contained in the spectrum of A, and contains the boundary of the spectrum of A.
\item
  Let \(f \in C[X]\) a polynomial. The image under f of the singular spectrum of a is the singular spectrum of \(f(a)\) .
\end{enumerate}

\begin{enumerate}
\def\labelenumi{\arabic{enumi})}
\setcounter{enumi}{16}
\tightlist
\item
  Let A be a normed algebra. For every \(x \in A\) , we set:
\end{enumerate}

\[
\lambda (x) = \inf _ {y \neq 0} \frac {\| x y \|}{\| y \|} \quad \lambda^ {\prime} (x) = \inf _ {y \neq 0} \frac {\| y x \|}{\| y \|}.
\]

\begin{enumerate}
\def\labelenumi{\alph{enumi})}
\tightlist
\item
  We have
\end{enumerate}

\[
| \lambda (x) - \lambda (y) | \leqslant \| x - y \|, \quad | \lambda^ {\prime} (x) - \lambda^ {\prime} (y) | \leqslant \| x - y \|,
\]

\[
\lambda (x) \lambda (y) \leqslant \lambda (x y) \leqslant \| x \| \lambda (y), \quad \lambda^ {\prime} (x) \lambda^ {\prime} (y) \leqslant \lambda^ {\prime} (x y) \leqslant \lambda^ {\prime} (x) \| y \|.
\]

\begin{enumerate}
\def\labelenumi{\alph{enumi})}
\setcounter{enumi}{1}
\tightlist
\item
  The set of left (resp. right) topological zero divisors in A is closed.
\end{enumerate}

\begin{enumerate}
\def\labelenumi{\arabic{enumi})}
\setcounter{enumi}{17}
\item
  Let A be a unital Banach algebra. The set of \(x\) which are neither invertible nor topological zero divisors, is open in A.
\item
  Let X be a complex Banach space and \(x \in A = \mathcal{L}(X)\) .
\end{enumerate}

\begin{enumerate}
\def\labelenumi{\alph{enumi})}
\tightlist
\item
  If x is not invertible, then x is a left or right topological zero divisor. (Consider successively the following cases:)
\end{enumerate}

\(1^{\circ}\) ) \(x\) is not injective;

\(2^{\circ})\overline{x(\mathrm{X})}\neq \mathrm{X};\)

\(3^{\circ}) x \text{ is injective, } x(\mathrm{X}) \text{ is dense in X and distinct from X.}\)

\begin{enumerate}
\def\labelenumi{\alph{enumi})}
\setcounter{enumi}{1}
\tightlist
\item
  If x maps X bicontinuously onto a closed vector subspace of X distinct from X, or if x is non-injective with image X, then x is interior to the set of non-invertible elements of A.
\end{enumerate}

\begin{enumerate}
\def\labelenumi{\arabic{enumi})}
\setcounter{enumi}{19}
\item
  In the algebra A of Example 9 of I, p.~20, the identity function z is neither invertible nor a topological zero divisor.
\item
  Let A be a unital normed algebra. Let B be the normed algebra \(\mathcal{L}(\mathrm{A})\) . Then the image of A under the morphism \(x \mapsto \gamma_x\) is a full subalgebra of B. Therefore, if \(x \in \mathrm{A}\) , one has \(\mathrm{Sp}_{\mathrm{A}}(x) = \mathrm{Sp}_{\mathrm{B}}(\gamma_x)\) .
\end{enumerate}

\begin{enumerate}
\def\labelenumi{¶ \arabic{enumi})}
\setcounter{enumi}{21}
\tightlist
\item
  Let A be a Banach algebra.
\end{enumerate}

\begin{enumerate}
\def\labelenumi{\alph{enumi})}
\tightlist
\item
  Let S be the set of bounded sequences of elements of A, with the norm \(\|(\boldsymbol{x}_{n})\| = \sup \|x_{n}\|\) . Let R be the set of \((x_{n}) \in \mathrm{S}\) such that \(\|x_{n}\|\) tends to 0. Then S is a Banach algebra and R is a closed two-sided ideal of S.
\end{enumerate}

Let \(S' = \text{S/R}\) and \(\varphi: \text{S} \to \text{S}'\) the canonical homomorphism.

\begin{enumerate}
\def\labelenumi{\alph{enumi})}
\setcounter{enumi}{1}
\item
  The map \(\theta\) from A into S' defined by \(\theta(x) = \varphi(x, x, x, \ldots)\) is an isometric isomorphism of A onto a subalgebra of S'. Every left topological divisor of zero in S' is a left divisor of zero. For \(x \in A\) to be a left topological divisor of zero, it is necessary and sufficient that \(\theta(x)\) be a left divisor of zero in S'.
\item
  Assume that A has a unit element. Prove that there exists a Banach algebra B and an isometric isomorphism of A onto a subalgebra of B such that an element of A is non-invertible if and only if its image in B is a left or right divisor of zero. (Use (a) and Exercises 19 and 21.)
\end{enumerate}

\begin{enumerate}
\def\labelenumi{\arabic{enumi})}
\setcounter{enumi}{22}
\tightlist
\item
  Let A be a normed algebra and R its radical.
\end{enumerate}

\begin{enumerate}
\def\labelenumi{\alph{enumi})}
\item
  If \(x \in R\) , then x is quasi-nilpotent. (We have \(\mathrm{Sp}_{\mathrm{A}}^{\prime}(x) = \{0\}\) and a fortiori \(\mathrm{Sp}_{\widehat{\mathrm{A}}}^{\prime}(x) = \{0\}\) .)
\item
  Suppose that A is a Banach algebra. Let I be a left ideal of A all of whose elements are quasi-nilpotent. Then \(I \subset R\) (We may assume that A has a unit element. If \(x \in R\) and \(a \in A\) , one has \(\mathrm{Sp}_{\mathrm{A}}(ax) = \{0\}\) , hence 1 - ax is invertible; hence \(x \in R\) .
\end{enumerate}

\begin{enumerate}
\def\labelenumi{\arabic{enumi})}
\setcounter{enumi}{23}
\item
  Let A be a complex Banach algebra, B a complex algebra without radical, and \(\varphi\) a morphism from A onto B. The kernel of \(\varphi\) is closed.
\item
  Let A be a complex Banach algebra and \(\pi\) a nonzero irreducible representation of A in a complex vector space X. Let \(\xi_{0}\) be a nonzero element of X.
\end{enumerate}

\begin{enumerate}
\def\labelenumi{\alph{enumi})}
\item
  The annihilator I of \(\xi_0\) is a maximal regular left ideal of A; it is closed.
\item
  The map \(x \mapsto x\xi_{0}\) defines, by passage to the quotient, an isomorphism \(\varphi\) of the A-module A/I onto the A-module X.
\item
  If one transports by \(\varphi\) the norm of the Banach space A/I, then X becomes a Banach space, and \(\| \pi (x)\| \leqslant \| x\|\) for all \(x\in \mathrm{A}\) .
\item
  If A is primitive (that is, \(\{0\}\) is a primitive ideal, cf. definition I, p. 11, cf. also A, VIII, p. 433, exerc. 6), then \(A = \{0\}\) or \(C \cdot 1\) according as A does or does not have a unit element. (Use the preceding and cor. 4 of I, p. 26).
\end{enumerate}

\begin{enumerate}
\def\labelenumi{¶ \arabic{enumi})}
\setcounter{enumi}{25}
\tightlist
\item
  Let A be a Banach algebra and R its radical.
\end{enumerate}

\begin{enumerate}
\def\labelenumi{\alph{enumi})}
\tightlist
\item
  Prove that if \(r \in R\) , the series
\end{enumerate}

\[
- \frac {1}{2} \sum_ {k = 1} ^ {\infty} \binom{1 / 2}{k} (- 4 r) ^ {k}
\]

converges to an element \(x \in R\) such that \(x^{2} - x + r = 0\) .

\begin{enumerate}
\def\labelenumi{\alph{enumi})}
\setcounter{enumi}{1}
\tightlist
\item
  Let u be an element of A whose class in A/R is idempotent. There exists an idempotent of A congruent to u modulo R. (We may assume that A has a unit element. Let \(q = u - u^{2} \in R\) and \(x \in R\) a solution, constructed with the aid of a), of \(x^{2} - x - q(1 - 4q)^{-1} = 0\) . Then \(u - (2u - 1)x\) answers the question.)
\end{enumerate}

\begin{enumerate}
\def\labelenumi{\arabic{enumi})}
\setcounter{enumi}{26}
\item
  Let A be a complex unital Banach algebra, \(x \in A\) and \(\zeta\) a boundary point of \(\mathrm{Sp}_{\mathrm{A}}(x)\) . The resolvent of x cannot be extended to a continuous function at \(\zeta\) .
\item
  Let A be a complex unital Banach algebra, \(\Delta\) an open subset of C and \(\lambda \mapsto \mathrm{R}(\lambda)\) a map of \(\Delta\) to A such that
\end{enumerate}

\[
\mathrm{R} (\lambda) - \mathrm{R} (\mu) = - (\lambda - \mu) \mathrm{R} (\lambda) \mathrm{R} (\mu)\tag{5}
\]

for all \(\lambda,\mu\in\Delta\) .

\begin{enumerate}
\def\labelenumi{\alph{enumi})}
\tightlist
\item
  Let \(\lambda_0\in \Delta\) . For every \(\lambda \in \Delta\) such that \(|\lambda -\lambda_0|\cdot \| \mathrm{R}(\lambda)\| <  1\) , we have:
\end{enumerate}

\[
\mathrm{R} (\lambda) = \sum_ {n = 0} ^ {\infty} (\lambda_ {0} - \lambda) ^ {n} \mathrm{R} (\lambda_ {0}) ^ {n + 1}.
\]

\begin{enumerate}
\def\labelenumi{\alph{enumi})}
\setcounter{enumi}{1}
\tightlist
\item
  For every \(\lambda \in \Delta\) and every integer \(n\geqslant 0\) , one has
\end{enumerate}

\[
\frac {\partial^ {n}}{\partial \lambda^ {n}} \mathrm{R} (\lambda) = (- 1) ^ {n} n! \mathrm{R} (\lambda) ^ {n + 1}.
\]

\begin{enumerate}
\def\labelenumi{\alph{enumi})}
\setcounter{enumi}{2}
\tightlist
\item
  If R is holomorphic at infinity, then there exists \(z, j, x \in \mathrm{A}\) such that \(z^2 = 0\) , \(j^2 = j\) , \(zj = jz = 0\) , \(x \in j\mathrm{Aj}\) , and \(\mathrm{R}(\lambda) = z + j\mathrm{R}(\lambda, x)\) for \(|\lambda|\) sufficiently large. (Write \(\mathrm{R}(\lambda) = \sum_{n=0}^{\infty} c_n \lambda^{-n}\) for \(|\lambda| > \lambda_0\) , and express that R satisfies (5). One may set \(c_0 = z\) , \(c_1 = j\) , \(c_2 = x\) .)
\end{enumerate}

\begin{enumerate}
\def\labelenumi{\alph{enumi})}
\setcounter{enumi}{3}
\tightlist
\item
  For there to exist \(a \in A\) such that \(R\) is the restriction to \(\Delta\) of the resolvent \(\lambda \mapsto R(a, \lambda)\) of \(a\) , it is necessary and sufficient that there exist \(\lambda_0 \in \Delta\) such that \(R(\lambda_0)\) be invertible.
\end{enumerate}

\begin{enumerate}
\def\labelenumi{\alph{enumi})}
\setcounter{enumi}{4}
\tightlist
\item
  For every \(x \in A\) , the function \(\lambda \mapsto \sum_{n=0}^{\infty} (\lambda_0 - \lambda)^n x^{n+1}\) satisfies (5) for \(|\lambda - \lambda_0| \cdot \|x\| < 1\) .
\end{enumerate}

\begin{enumerate}
\def\labelenumi{\arabic{enumi})}
\setcounter{enumi}{28}
\tightlist
\item
  \begin{enumerate}
  \def\labelenumii{\alph{enumii})}
  \tightlist
  \item
    Let E be a complex Banach space, \(f: \mathbf{C} \to \mathbf{E}\) an entire function and \(\theta \mapsto \mathrm{P}(\theta) = \sum_{\mu=k}^{m} c_{\mu} e^{\mu i\theta}\) a trigonometric polynomial ( \(c_k \neq 0\) ). We assume that \(|\mathrm{P}(\theta)| \cdot \|f(re^{i\theta})\| \leqslant \mathrm{Mr}^{\alpha}\) for \(r \geqslant r_0\) (M, \(r_0\) , \(\alpha\) being constants \(\geqslant 0\) ). Then \(f\) is a polynomial of degree \(\leqslant \alpha\) . (Let \(f(\zeta) = \sum_{\nu=0}^{\infty} a_{\nu} \zeta^{\nu}\) , and
  \end{enumerate}
\end{enumerate}

\[
\mathrm{J} (r, p) = r ^ {- p} \int_ {0} ^ {2 \pi} \mathrm{P} (\theta) f (r e ^ {i \theta}) e ^ {- (k + p) i \theta} d \theta .
\]

Compute the limit of \(\mathrm{J}(r,p)\) when \(r\to+\infty\) in two different ways for \(p>\alpha.\)

\begin{enumerate}
\def\labelenumi{\alph{enumi})}
\setcounter{enumi}{1}
\tightlist
\item
  With the notations of \(a\) ), if \(\alpha\) is an integer and if, for every \(\theta \in \mathbf{R}\) , one has
\end{enumerate}

\[
\lim _ {r \to + \infty} r ^ {- \alpha} | \mathrm{P} (\theta) | \| f (r e ^ {i \theta}) = 0,
\]

then \(f\) is a polynomial of degree \(< \alpha\) .

\begin{enumerate}
\def\labelenumi{\alph{enumi})}
\setcounter{enumi}{2}
\tightlist
\item
  Let A be a unital complex Banach algebra and x a quasi-nilpotent element of A. We assume that, for some integer n \textgreater{} 0, and for every \(\theta \in R\) , one has
\end{enumerate}

\[
\lim _ {r \to + \infty} r ^ {- n} | \cos^ {n} \theta | \cdot \| (1 - r e ^ {i \theta} x) ^ {- 1} \| = 0.
\]

Then \(x^n = 0\) . (Use \(b\) ).)

\begin{enumerate}
\def\labelenumi{¶ \arabic{enumi})}
\setcounter{enumi}{29}
\tightlist
\item
  a) Let E be a complex Banach space and \(x: \mathbf{C}=\{1\}\to\mathrm{E}\) an entire function of \(1/(\zeta-1)\) . One denotes by \(x(\zeta)=\sum_{n=0}^{\infty}a_{n}\zeta^{n}\) , \(x(\zeta)=\sum_{n=0}^{\infty}b_{n}\zeta^{-n}\) the expansions of \(x\) for \(|\zeta|<1\) , \(|\zeta|>1\) respectively. We assume that there exists \(\alpha>0\) such that \(\|a_{n}\|=o(n^{\alpha})\) , \(\|b_{n}\|=o(n^{\alpha})\) . Then \(x(\zeta)\) is a polynomial in \(1 / (\zeta - 1)\) of degree \(< \alpha + 1\) . (Let \(\varepsilon > 0\) . We have \(\|x(\zeta)\| \leqslant \varepsilon(1 - r)^{-1 - \alpha}\) for \(r_{\varepsilon} \leqslant r = |\zeta| < 1\) , \(\|x(\zeta)\| \leqslant \varepsilon(1 - r^{-1})^{-1 - \alpha}\) for \(1 < r \leqslant r_{\varepsilon}'\) . Setting \(1 - \zeta = (\omega + \frac{1}{2})^{-1}\) , \(\omega = \mathrm{Re}^{i\theta}\) , show that, for \(0 \leqslant r < 1\) and \(\mathrm{R} \geqslant 1\) , one has
\end{enumerate}

\[
(1 - r) ^ {- 1} \leqslant \frac {9}{4} \frac {\mathrm{R}}{\cos \theta}
\]

and that, for 1 \textless{} r and R ≥ 1,

\[
(1 - r ^ {- 1}) ^ {- 1} \leqslant \frac {9}{4} \frac {\mathrm{R}}{\cos \theta}.
\]

Set \(y(\omega) = x(\zeta)\) and show that \(R^{-\alpha - 1}|\cos \theta|^{\alpha + 1}\| y(Re^{i\theta})\|\) tends to 0 as \(R\) tends to \(+\infty\) , uniformly in \(\theta\) . Then apply exerc. 29, b).

Let A be a complex unital Banach algebra, \(q \in A\) a quasi-nilpotent element and \(x = 1 + q\) .

\begin{enumerate}
\def\labelenumi{\alph{enumi})}
\setcounter{enumi}{1}
\item
  For \(q^{N}=0\) , it is necessary and sufficient that \(\|x^{\pm n}\|=o(n^{N})\) when n tends to \(+\infty\) . (To see that the condition is sufficient, show, by applying a), that \(\mathrm{R}(\lambda,x)\) is a polynomial in \(1/(\lambda-1)\) of degree \(\leqslant N\) . On the other hand, \(\mathrm{R}(\lambda,x)=\sum_{n=0}^{\infty}q^{n}(\lambda-1)^{-n-1}\) .)
\item
  Let R be the radical of A and G a subgroup of the group of invertible elements of A. If, for every \(x \in G\) , one has \(\|x^{n}\| = o(n)\) , \(\|x^{-n}\| = o(n)\) for \(n \to +\infty\) , then the restriction to G of the canonical map \(A \to A/R\) is injective.
\end{enumerate}

\begin{enumerate}
\def\labelenumi{\arabic{enumi})}
\setcounter{enumi}{30}
\tightlist
\item
  \begin{enumerate}
  \def\labelenumii{\alph{enumii})}
  \tightlist
  \item
    Let A be a normed algebra, and x, y two elements of A such that xy - yx = y. Prove that \(y^{n} = 0\) for \(n > 2\|x\|\) . (Use the formula \(xy^{n} - y^{n}x = ny^{n}\) ).
  \end{enumerate}
\end{enumerate}

\begin{enumerate}
\def\labelenumi{\alph{enumi})}
\setcounter{enumi}{1}
\item
  Let L be a finite-dimensional Lie algebra over C. Prove that, if L is not nilpotent, it contains two nonzero elements x, y such that \([x, y] = y\) . (Use the fact that there exists an \(x \in L\) such that \(\operatorname{ad}(x)\) is not nilpotent.)
\item
  Let L be a real or complex finite-dimensional Lie algebra such that the enveloping algebra of L possesses a norm compatible with its algebra structure. Prove that L is nilpotent. (Use a) and b).
\end{enumerate}

\begin{enumerate}
\def\labelenumi{\alph{enumi})}
\setcounter{enumi}{3}
\tightlist
\item
  Let V be a finite-dimensional real or complex vector space, \(\mathfrak{L}\) a Lie subalgebra of \(\mathfrak{gl}(\mathrm{V})\) consisting of nilpotent endomorphisms, and \(x_{1},\ldots ,x_{n}\) elements of \(\mathfrak{L}\) . Let \(\alpha >0\) . Prove that there exists a Hilbert space structure on V such that \(\| x_i\| \leqslant \alpha\) for \(1\leqslant i\leqslant n\) .
\end{enumerate}

\begin{enumerate}
\def\labelenumi{\alph{enumi})}
\setcounter{enumi}{4}
\tightlist
\item
  Let \(\mathfrak{L}\) be a real or complex finite-dimensional nilpotent Lie algebra, and U its enveloping algebra. Prove that there exists on U a norm compatible with its algebra structure. (Using the method of LIE, I, §3, exerc. 5 and §7, exerc. 3, show that there exists a family \((\pi_{\lambda})\) of representations of \(\mathfrak{L}\) in finite-dimensional spaces \(V_{\lambda}\) , having the following property:
\end{enumerate}

for all \(u \in U\) , there exists an \(\lambda\) such that \(\pi_{\lambda}(u) \neq 0\) , and \(\pi_{\lambda}(\mathfrak{L})\) consists, for every \(\lambda\) , of nilpotent endomorphisms. By endowing each \(V_{\lambda}\) with a Hilbert structure furnished by d), one obtains in the Hilbert space V, the Hilbert sum of the \(V_{\lambda}\) , a representation \(\pi\) of L by continuous endomorphisms, whence an injective morphism \(\varphi\) of U into \(\mathcal{L}(V)\) . Set \(\|u\|_{U} = \|\varphi(u)\|\) for all \(u \in U\) ).

\begin{enumerate}
\def\labelenumi{\alph{enumi})}
\setcounter{enumi}{5}
\tightlist
\item
  Let A be a complex unital normed algebra and x and y elements of A. If \(A \neq \{0\}\) , one has xy - yx ≠ 1. (First proof: one has \(\mathrm{Sp}_{\mathrm{A}}^{\prime}(xy) = \mathrm{Sp}_{\mathrm{A}}^{\prime}(yx)\) ; if xy = yx + 1, then \(\mathrm{Sp}_{\mathrm{A}}(xy)\) is deduced from \(\mathrm{Sp}_{\mathrm{A}}(yx)\) by the translation \(z \mapsto z + 1\) ; deduce that \(\mathrm{Sp}_{\mathrm{A}}(xy)\) is unbounded, which is absurd. Second proof: if \([x, y] = 1\) , one has \([xy, y] = y\) , whence \(y^{n} = 0\) for n sufficiently large by a); since \([x, y^{p}] = py^{p-1}\) deduce that y = 0).
\end{enumerate}

\begin{enumerate}
\def\labelenumi{\arabic{enumi})}
\setcounter{enumi}{31}
\item
  Let A be a normed algebra and \(x \in A\) . One says that x is topologically nilpotent if \(x^{n}\) tends to 0 when n tends to \(+\infty\) . Show that x is topologically nilpotent if and only if \(\varrho(x) < 1\) .
\item
  Let A be a commutative unital Banach algebra over C. Suppose that every closed ideal I of A is finitely generated (that is, there exist \(n \in N\) and \(x_{1}, \ldots, x_{n} \in A\) such that \(I = Ax_{1} + \cdots + Ax_{n}\) ).
\end{enumerate}

\begin{enumerate}
\def\labelenumi{\alph{enumi})}
\tightlist
\item
  Every ideal I of A is closed. (Write \(\bar{\mathrm{I}} = \mathrm{A}x_{1} + \dots +\mathrm{A}x_{n}\) with \(x_{1},\ldots ,x_{n}\) in A. Let \((x_{ip})_{p\geqslant 1}\) be a sequence of elements of I tending to \(x_{i}\) . For every \((y_{1},\ldots ,y_{n})\in \mathrm{A}^{n}\) and every \(p\geqslant 1\) , set
\end{enumerate}

\[
\begin{array}{c} \varphi (y _ {1}, \ldots , y _ {n}) = x _ {1} y _ {1} + \dots + x _ {n} y _ {n} \in \bar {\mathrm{I}} \\ \varphi_ {p} (y _ {1}, \ldots , y _ {n}) = x _ {i p} y _ {1} + \dots + x _ {n p} y _ {n} \in \mathrm{I}. \end{array}
\]

We have \(\varphi\in\mathcal{L}(A^{n},\bar{I})\) , \(\varphi_{p}\in\mathcal{L}(A^{n},\bar{I})\) , \(\|\varphi-\varphi_{p}\|\) tends to 0 as \(p\) tends to \(+\infty\) , and \(\varphi\) is surjective. Deduce that \(\varphi_{p}\) is surjective for \(p\) large enough, considering the transposes of \(\varphi\) and \(\varphi_{p}\) Hence I= \(\bar{I}\) .

\begin{enumerate}
\def\labelenumi{\alph{enumi})}
\setcounter{enumi}{1}
\item
  If A is an integral domain, \(A = C \cdot 1\) (Use a) and Exerc. 15.)
\item
  If the only nilpotent element of A is 0, there exists an integer \(n \geqslant 0\) such that A is isomorphic to \(\mathbf{C}^n\) . (Since A is Noetherian, \(\{0\}\) is the intersection of prime ideals \(\mathfrak{p}_1, \ldots, \mathfrak{p}_m\) ; these are closed by \(a\) ), and the algebras \(\mathrm{A}/\mathfrak{p}_i\) are isomorphic to \(\mathbf{C}\) according to \(b\) .)
\end{enumerate}

\begin{enumerate}
\def\labelenumi{\alph{enumi})}
\setcounter{enumi}{3}
\tightlist
\item
  We have \(\dim_{\mathbf{C}}\mathrm{A} < +\infty\) . (Let N be the set of nilpotent elements of A, which is a (closed) ideal of A. Then \(\dim (\mathrm{A / N}) < +\infty\) according to \(c\) ). Since N is a finitely generated ideal, \(\mathrm{N}^n = 0\) for \(n\) sufficiently large. Finally, each \(\mathrm{N}^i /\mathrm{N}^{i + 1}\) is a finitely generated module over A/N, hence \(\dim (\mathrm{N}^i /\mathrm{N}^{i + 1}) < +\infty\) .)
\end{enumerate}

\begin{enumerate}
\def\labelenumi{\arabic{enumi})}
\setcounter{enumi}{33}
\tightlist
\item
  Let A be a unital algebra over C, equipped with a separated locally convex topology such that multiplication in A is separately continuous.
\end{enumerate}

. An element \(a \in A\) is said to be regular if there exists \(r \geqslant 0\) such that \(a - \lambda\) let be invertible for \(|\lambda| \geqslant r\) , and such that the set of \((a - \lambda)^{-1}\) , for \(|\lambda| \geqslant r\) , be bounded.

\begin{enumerate}
\def\labelenumi{\alph{enumi})}
\item
  If a is regular, \(\lambda(a-\lambda)^{-1}\) remains bounded for \(|\lambda|\geqslant r\) .
\item
  If the mapping \(x \mapsto x^{-1}\) is defined in a neighborhood of 1 and continuous at the point 1, every element of A is regular.
\item
  For every \(a \in A\) , let \(U_{a}\) the set of \(\lambda \in C\) such that \((a - \mu)^{-1}\) exists and is bounded for all \(\mu\) of a neighborhood of \(\lambda\) . Let \(S_{a} = C - U_{a}\) . Then \(S_{a}\) is closed; and, if a is regular, \(S_{a}\) is compact nonempty (provided that \(A \neq \{0\}\) ).
\item
  Let \(a \in A\) . The function \(\lambda \mapsto (a - \lambda)^{-1}\) is holomorphic in \(U_{a}\) .
\item
  Let \(a \in A\) . For \(\lambda \in U_{a}\) , it is necessary and sufficient that \(a - \lambda\) has a regular inverse.
\item
  If every element of A is regular, and if A is a field, then we have \(A = C \cdot 1\) .
\end{enumerate}

\begin{enumerate}
\def\labelenumi{\arabic{enumi})}
\setcounter{enumi}{34}
\tightlist
\item
  Let A be an algebra over R which is a field, equipped with a separated locally convex topology such that multiplication in A is continuous, and such that the map \(x \mapsto x^{-1}\) be continuous at 1. Then A is R-isomorphic either to R, to C, or to the quaternion field H. (If A is commutative and there exists \(u \in A\) with \(u^{2} = -1\) give A a structure of algebra over C, and apply exerc. 34, f).) If A is commutative and there exists no element \(u \in A\) such that \(u^{2} = -1\) apply exerc. 34, f) to the field \(A \otimes_{R} C\) If A is noncommutative, argue as in AC, VI, §6, no. 4, th. 1, third case.)
\end{enumerate}

\begin{enumerate}
\def\labelenumi{¶ \arabic{enumi})}
\setcounter{enumi}{35}
\tightlist
\item
  Let A be the algebra over C consisting of the restrictions to {[}0,1{]} of rational functions in one variable with complex coefficients. This algebra is a field.
\end{enumerate}

\begin{enumerate}
\def\labelenumi{\alph{enumi})}
\item
  We endow A with the topology of convergence in measure (INT, IV, §5, no. 11). Then the mappings \((x,y)\mapsto x+y\) and \((x,y)\mapsto xy\) of A × A into A are continuous, and the map \(x\mapsto x^{-1}\) from A* to A is continuous. The topology of A is not locally convex.
\item
  For \(n \in N^{*}\) , we define the sequence \((w_{r,n})_{r \in \mathbf{Z}}\) by \(w_{r,n} = (-r + 1)^{n(-r+1)}\) if r \textless{} 0, \(w_{0,n} = 1\) and \(w_{r,n} = (r + 1)^{-(r+1)/n}\) if r \textgreater{} 0. If \(f \in A\) , set \(p_n(f) = \sum_r w_{r,n} |a_r| < +\infty\) , where \(\sum a_r t^r\) is the Laurent expansion of \(f(t)\) at 0. Then the \(p_n\) are semi-norms that define on A a metrizable locally convex topology. The map \((x, y) \mapsto xy\) The mapping from A × A to A is continuous. The map \(x \mapsto x^{-1}\) The map from A* into A is not continuous. The algebra A is not complete for this topology.
\end{enumerate}

\begin{enumerate}
\def\labelenumi{\arabic{enumi})}
\setcounter{enumi}{36}
\tightlist
\item
  \begin{enumerate}
  \def\labelenumii{\alph{enumii})}
  \tightlist
  \item
    Let A be an algebra over C equipped with a locally convex topology. The following conditions are equivalent:
  \end{enumerate}
\end{enumerate}

\begin{enumerate}
\def\labelenumi{(\roman{enumi})}
\item
  there exists a fundamental system \((\mathrm{U}_{i})\) of convex balanced neighborhoods of 0 such that \(U_{i}U_{i} \subset U_{i}\) for all i ;
\item
  A is isomorphic to a subalgebra of a product of normed algebras.
\item
  the topology of A can be defined by a family of semi-norms \(p_{i}\) satisfying \(p_{i}(xy) \leqslant p_{i}(x)p_{i}(y)\) for all \(x, y \in A\) .
\end{enumerate}

If A satisfies these conditions, we say that A is locally m-convex.

\begin{enumerate}
\def\labelenumi{\alph{enumi})}
\setcounter{enumi}{1}
\item
  Let A be a locally m-convex algebra. The map \((x,y)\mapsto xy\) The mapping from A × A into A is continuous. If A is unital, and if G denotes the set of invertible elements of A, then the map \(x\mapsto x^{-1}\) of G into A is continuous.
\item
  Let A be a unital locally m-convex algebra. The spectrum of every element of A is nonempty. If A is a field, then \(A = C \cdot 1\) .
\item
  Let A be a complete locally m-convex algebra. There exists a directed set I, a family \((\mathrm{A}_{i})_{i\in\mathrm{I}}\) of Banach algebras, continuous morphisms \(\pi_{ij}\colon A_{j}\mapsto A_{i}\) for \(i\leqslant j\) , with \(\pi_{ij}\circ\pi_{jk}=\pi_{ik}\) such that A is isomorphic to the subalgebra of \(\prod_{i}A_{i}\) consisting of \((x_{i})\) such that \(\pi_{ij}(x_{j})=x_{i}\) for \(i\leqslant j\) .
\item
  The algebra \(\mathcal{C}_{\mathbf{C}}(\mathbf{R})\) , equipped with the topology of compact convergence, is locally m-convex. The set of invertible elements of \(\mathcal{C}_{\mathbf{C}}(\mathbf{R})\) is dense in \(\mathcal{C}_{\mathbf{C}}(\mathbf{R})\) ; it is not open.
\item
  Let D be an open subset of C, and let A be the algebra of complex-valued holomorphic functions on D, equipped with the topology of compact convergence. Then A is locally m-convex. The set of invertible elements of A is closed.
\item
  The algebra of infinitely differentiable complex-valued functions on {[}0,1{]}, equipped with the topology of uniform convergence of each derivative, is locally m-convex.
\item
  Let A be the algebra of (classes of) complex functions f on {[}0,1{]} such that \(f \in \mathrm{L}^{p}([0,1])\) for every p \textgreater{} 1. Endowed with the topology defined by the family of seminorms \(f \mapsto \|f\|_{p}\) (p \textgreater{} 1), A is not locally m-convex. The map \((x,y) \mapsto xy\) from A × A into A is continuous.
\end{enumerate}

